% ===================================================================
%  TAULA N° 2 — Lo Còrs Uman. — La Recreacion. — Los Jòcs.
%  Traduction 2026 d'apres texto/io, controlee sur texto/fr et
%  texto/ca. Consigne : voir texto/oc/10-taula-01.tex.
% ===================================================================

%%K t02-apar-1 apar
\VUsaut{4.0mm}
\VUtitre{15.0pt}{60}{TAULA N° 2}
\VUsaut{2.0mm}
\VUtitre{11.4pt}{0}{Lo Còrs Uman. --- La Recreacion.}
\VUtitre{11.4pt}{0}{Los Jòcs.}

%%K t02-tit-0 sub
\VUtitre{13.2pt}{0}{Lo Còrs Uman.}
\VUsaut{2.4mm}
\VUcentre{10.2pt}{0}{\textit{(Leiçon simpla de sciéncias naturalas.)}}
\VUsaut{3.0mm}

%%K t02-tit-1 sub pos=t02-01-1
\VUpk{10.2pt}{40}{Lo Cap.}
\VUsaut{3.0mm}

%%K t02-01-1 p
1. --- Las principalas partidas del còrs uman son: lo \VUgras{cap}, lo
\VUgras{tronc} e los \VUgras{membres}.

%%K t02-02-1 p
2. --- Lo cap compren doas partidas: lo \VUgras{cran}, que conten lo
\VUgras{cervèl} e que los \VUgras{pels} cobrisson, e la \VUgras{cara}.

%%K t02-03-1 p
3. --- Dins nòstre dessenh vesèm ni lo cran ni lo cervèl; mas vesèm amb tota
clartat las partidas importantas de la cara.

%%K t02-04-1 p
4. --- Vaicí d'en primièr lo \VUgras{front}, qu'anóncia una intelligéncia
poderosa, un pensament totjorn desrevelhat. Las \VUgras{tempes} son fòrça
liuras. L'\VUgras{uèlh} es viu e plan dobèrt. Una \VUgras{cilha} espessa e de
\VUgras{parpèlas} longas lo protegisson.

%%K t02-05-1 p
5. --- Vaicí en mai lo \VUgras{nas} e las \VUgras{narras}. Lo nas nos sèrv per
sentir e per respirar. De cada costat del nas i a las \VUgras{gautas}, e mai luènh
las \VUgras{aurelhas}. La partida inferiora de la cara es formada per la
\VUgras{boca} e lo \VUgras{mentonh}.

%%K t02-06-1 p
6. --- Las vesèm pas gaire plan, que las \VUgras{mostachas} e las \VUgras{patilhas}
nos las escondon. Mas sabèm que la boca a las doas \VUgras{pòtas}, la
\VUgras{maissa superiora} e la \VUgras{maissa inferiora} amb sas \VUgras{dents},
e la \VUgras{lenga}. Amb la boca parlam e manjam. Amb la lenga tastam los
aliments.

%%K t02-07-1 p
7. --- L'uèlh, l'aurelha, lo nas e la boca son, amb los dets, los organs dels
cinc sens: la \VUgras{vista}, l'\VUgras{ausida}, l'\VUgras{odorat}, lo
\VUgras{gost} e lo \VUgras{tocar}.

%%K t02-08-1 p
8. --- Lo cap es doncas una de las partidas mai interessantas del còrs uman.

%%K t02-tit-2 sub
\VUsaut{4.4mm}
\VUpk{10.2pt}{40}{Lo Tronc.}
\VUsaut{3.0mm}

%%K t02-09-1 p
9. --- Lo \VUgras{còl} unís lo cap al tronc. Compren la \VUgras{gòrja} e la
\VUgras{nuca}.

%%K t02-10-1 p
10. --- Lo tronc conten tanben fòrça organs essencials. Dins lo \VUgras{pièch} i a
los \VUgras{palmons}, amb los quals respiram, e lo \VUgras{còr}, qu'es lo centre
de la vida. Quand lo còr daissa de batre, l'èsser viu morís.

%%K t02-11-1 p
11. --- Las \VUgras{còstas} sostenon lo pièch.

%%K t02-12-1 p
12. --- Las autras partidas del tronc son lo \VUgras{flanc}, l'\VUgras{esquina},
las \VUgras{espatlas} e lo \VUgras{ventre}. Nos jasèm sul flanc o sus l'esquina
per dormir, e portam las cargas sus l'espatla.

%%K t02-13-1 p
13. --- Dins lo ventre i a l'\VUgras{estomac} e los \VUgras{budèls}. Nos sèrvon
per digerir los aliments.

%%K t02-tit-3 sub
\VUsaut{4.4mm}
\VUpk{10.2pt}{40}{Los Membres.}
\VUsaut{3.0mm}

%%K t02-14-1 p
14. --- Avèm quatre \VUgras{membres}: dos \VUgras{braces} e doas \VUgras{cambas}.
Amb los braces prenèm, tenèm, levam e recampam los objèctes; amb las cambas nos
tenèm drechs, caminam e corrèm.

%%K t02-15-1 p
15. --- L'\VUgras{espatla} unís lo braç al còrs. Un \VUgras{muscle} fòrça robust,
lo biceps, mòu lo braç, que lo \VUgras{coide} unís a l'\VUgras{avantbraç}. Lo
\VUgras{ponhet} forma l'articulacion de la \VUgras{man}, qu'acaba pels
\VUgras{dets} e las \VUgras{onglas}. Dins la man destriam una \VUgras{vena} (e una
artèria) per ont cor lo \VUgras{sang}.

%%K t02-16-1 p
16. --- Las partidas de la camba son: la \VUgras{cuèissa}, lo \VUgras{genolh} e lo
\VUgras{garret}, lo \VUgras{gras de la camba}, que la \VUgras{carn} n'es cobèrta
per la \VUgras{pèl}, la \VUgras{cavilha} e lo \VUgras{pè}. Los \VUgras{òsses} de
la camba son fòrça gròsses e fòrça fòrts. A la cima del pè i a los \VUgras{dets
del pè}. Las autras partidas son la \VUgras{planta} e lo \VUgras{talon}.

%%K t02-17-1 p
17. --- Tala es la descripcion gaireben completa del còrs uman.

%%K t02-17-2 p
\textit{Nòta.} --- \textit{Amb l'ajuda de l'expression «me fa mal... (lo cap,
etc.)», lo professor poirà repassar tot aqueste vocabulari del ponch de vista del
bon o del marrit} \VUgras{estat del còrs}.

%%K t02-tit-4 sub
\VUsaut{5.0mm}
\VUpk{10.2pt}{40}{Gimnastica.}
\VUsaut{2.4mm}
\VUcentre{10.2pt}{0}{\textit{(Contat per un dels escolans.)}}
\VUsaut{3.0mm}

%%K t02-18-1 p
18. --- Per èsser soples, agils e sans, devèm exercir las cambas e los braces. Per
aquò jogam e fasèm de \VUgras{gimnastica}.

%%K t02-19-1 p
19. --- Dins la setmana, aqueles dos exercicis se fan dins la cort de l'institut
(vejatz la taula n° 1). Mas lo dijòus e lo dimenge anam dins un grand
\VUgras{prat}, ont avèm plaça per córrer e nos divertir.

%%K t02-20-1 p
20. --- Nòstres mèstres e nòstres parents nos i acompanhan de còps; mas es lo
\VUgras{professor de gimnastica} \textsuperscript{(12)} que dirigís los exercicis.

%%K t02-21-1 p
21. --- Dins lo prat, a esquèrra de l'intrada, i a los \VUgras{aparelhs de
gimnastica} \textsuperscript{(1)}: nos enfilam per l'\VUgras{escala}
\textsuperscript{(2)}, fasèm lo pin dins los \VUgras{anèls}
\textsuperscript{(3)} e dins lo \VUgras{trapèzi} \textsuperscript{(4)}, nos
issam amb las mans fins a la cima de l'\VUgras{escala de còrda}
\textsuperscript{(5)} o de la \VUgras{còrda de nosèls} \textsuperscript{(6)}.

%%K t02-22-1 p
22. --- Pendent aqueste temps, nòstres companhs sautan per dessús lo
\VUgras{caval de fusta} \textsuperscript{(7)} o las \VUgras{barras parallèlas}
\textsuperscript{(11)}. D'autres \VUgras{bòxan} \textsuperscript{(8)} o fan de
\VUgras{esgrima} \textsuperscript{(9)} amb masqueta e \VUgras{floret}. D'autres,
enfin, lèvan de \VUgras{pes} \textsuperscript{(10)}.

%%K t02-tit-5 sub
\VUsaut{4.6mm}
\VUpk{10.2pt}{40}{Los Jòcs.}
\VUsaut{3.0mm}

%%K t02-23-1 p
23. --- A l'entorn dels gimnastas, dròlles e dròllas jògan a de \VUgras{jòcs}
diferents. Las dròllas se divertisson amb son \VUgras{cèrcle}
\textsuperscript{(14)}; sautan a la \VUgras{còrda} \textsuperscript{(13)}, o
convidan sos fraires e sos cosins a una partida de \VUgras{cròquet}
\textsuperscript{(15)}. Las encanta d'escartar la \VUgras{bòla} de l'adversari amb
son \VUgras{maçòl de fusta}, just quand aqueste se pensava passar sens pena per
l'arquet!

%%K t02-24-1 p
24. --- Los dròlles aiman mai los jòcs mai rudes. Jògan de bon grat a las
\VUgras{quilhas} \textsuperscript{(16)}, que las abaton amb adreça amb la
\VUgras{bòla} \textsuperscript{(17)}. Mespresan pas tanpauc la \VUgras{baudufa}
\textsuperscript{(18)} e lo \VUgras{bilboquet} \textsuperscript{(19)}, ni
quitament lo \VUgras{jòc de la granolha} \textsuperscript{(20)}. Mas sos jòcs
preferits son las \VUgras{bilhas} \textsuperscript{(21)} e la \VUgras{pelòta a la
paret} \textsuperscript{(22)}, que la paret de la \VUgras{sèrra}
\textsuperscript{(26)} va fòrça plan per far rebombir la pelòta leugièra.

%%K t02-25-1 p
25. --- Lo pichon Joan, enfilat sus sas \VUgras{cambas de fusta}
\textsuperscript{(24)}, es fòrça adrech e sap gardar plan l'equilibri. A prèp
d'el, qualques companhs jògan a \VUgras{rescondedons} \textsuperscript{(25)} dins
aquela sèrra e darrièr la \VUgras{bartassa}. D'autres aiman mai la
\VUgras{man cauda} \textsuperscript{(28)} o lo \VUgras{volant}
\textsuperscript{(29)}, que volastreja d'una \VUgras{raqueta} a l'autra.

%%K t02-noto-1 noto
\VUnotes{143.2mm}{%
(*) Los jòcs d'enfants an sovent de noms purament nacionals. Los avèm nomenats en
ido tan plan coma avèm pogut, siá en nos inspirant de l'accion meteissa que los
caracteriza, siá en adoptant lo nom nacional d'un país o d'un autre quand es pas
tròp inexacte. Per exemple: lo \VUgras{jòc de la barrica} (alemand e francés) es
lo \VUgras{jòc de la granolha} \textsuperscript{(20)} (anglés), ont los jogaires
ensajan de getar sos discs redonds per la boca de la granolha de metal qu'es amb
la boca dobèrta sus la caissa (la barrica) traucada. Lo \VUgras{saut a l'espatla}
\textsuperscript{(30)} se destria del \VUgras{saut a l'esquina} solament en
aquò: per sautar per dessús lo cap, los jogaires apiejan las mans sus las
espatlas del companh, mentre que dins lo «\VUgras{saut a l'esquina}» las apiejan
solament sus son esquina. --- O ben, qualques uns dels que jògan s'inclinan
mentre que los autres sautan sus lors esquinas en apiejant las mans sus
l'espatla; los primièrs devon supportar lo còp sens se plegar, los segonds devon
sautar sens pèrdre l'equilibri (vejatz la taula meteissa).%
}

%%K t02-26-1 p
26. --- Al mitan del prat i a dos arbres. Al pè d'un d'eles, sèt o uèch dròlles
an organizat una partida de \VUgras{saut a l'espatla} \textsuperscript{(30)} (*).
Aqueste jòc se coneis pas dins totes los païses; mas tot lo mond sap çò qu'es lo
\VUgras{saut a l'esquina}, al qual los dròlles jògan pas ara.

%%K t02-tit-6 sub
\VUsaut{5.0mm}
\VUpk{10.2pt}{40}{Los Jòcs a l'aire liure.}
\VUsaut{3.4mm}

%%K t02-27-1 p
27. La \VUgras{galina bòrnia} \textsuperscript{(31)} es tanben fòrça divertenta;
dròllas e dròlles se pausan a torn de ròtle la \VUgras{benda} suls uèlhs e agafan
los maladrechs que se daissan prene.

%%K t02-28-1 p
28. --- Al mitan del prat, vaicí un jòc fòrça francés, quitament un pauc brusc:
es l'\VUgras{orsa} \textsuperscript{(32)}, fòrça mai bruchosa que lo jòc tan
tranquil de l'\VUgras{estèla} \textsuperscript{(33)} o que lo jòc anglés del
\VUgras{tenís} \textsuperscript{(34)}.

%%K t02-29-1 p
29. --- Aqueste darrièr espòrt fa las delícias dels grands escolans. Nòstres
professors eles meteisses lo practican de bon grat. Exigís fòrça adreça e
agilitat, e a ieu m'agrada fòrça. Daissi lo \VUgras{balançador}
\textsuperscript{(35)} per las dròllas.

%%K t02-30-1 p
30. --- M'agrada pas tanpauc un autre jòc anglés que pòt venir perilhós.

%%K t02-30-1b p
Parli del \VUgras{fotbòl} \textsuperscript{(36)}, que fa lo bonur de mon fraire
màger. Aimi mai nòstre vièlh \VUgras{jòc de barras} \textsuperscript{(38)}, tan
san e fòrça mens brutal.

%%K t02-tit-7 sub
\VUsaut{4.6mm}
\VUpk{10.2pt}{40}{Ciclisme e jòcs de pelòta.}
\VUsaut{3.0mm}

%%K t02-31-1 p
31. --- Tanben fau de bon grat un torn del prat en \VUgras{bicicleta}
\textsuperscript{(37)}.

%%K t02-32-1 p
32. --- L'an que ven aprendrai de \VUgras{montar a caval}
\textsuperscript{(39)}. Qu'es bèl un caval que camina o que tròta amb gràcia, e
que sauta los obstacles al galaup!

%%K t02-33-1 p
33. --- L'estiu aprenèm de \VUgras{nadar} \textsuperscript{(40)}. Nos capbussam e
nos banham amb delici dins l'aiga clara e fresca del riu. De còps, a la fin,
daissam anar un \VUgras{balon de papièr} \textsuperscript{(41)} que puja
majestuosament per l'aire tanlèu qu'es emplenat d'aire caud.

%%K t02-34-1 p
34. --- I a encara fòrça autres jòcs, mas acabariái pas jamai de los nomenar, e
per aqueste brèu resumit se vei que lo dijòus dins nòstre bèl prat son pas
brica enuèjoses.

%%K t02-34-2 p
N. B. --- \textit{Lo professor completarà aisidament aqueste capítol en descrivent
amb mai de detalh cadun dels jòcs importants.}
\VUsaut{6.0mm}
\VUfilet{22mm}
